\section{Introduction}\label{sec:intro}
\subsection{The problem that generalized thermoelasticity addresses}
Fourier's law makes the heat flux proportional to the temperature gradient at the same instant. Under ordinary conditions this is an excellent description of conduction in solids. Combined with energy balance it yields a parabolic equation, and that equation carries a property which is physically uncomfortable: a temperature disturbance applied at one point changes the temperature everywhere else at the same instant, however small the change may be.

For slow processes in bulk matter the difference between this instantaneous response and what a real, finite-speed carrier of heat would do is immeasurably small. It becomes measurable, and sometimes dominant, when the heating or loading is very fast (femtosecond to nanosecond laser pulses), when the structure is very small or very cold so that phonon mean free paths approach the structure size, when the medium is strongly heterogeneous, or when the observed quantity is a high-frequency elastic wave whose period is comparable to a thermal relaxation time.

Thermoelasticity couples the thermal field to deformation. Heating produces thermal stress through thermal expansion; deformation changes the temperature through the dilatation term in the energy balance. The classical coupled theory of Biot \citep{Biot1956} builds on the thermomechanical stress--strain law of Duhamel and Neumann \citep{Duhamel1837} and on the early dynamic solutions of Danilovskaya \citep{Danilovskaya1950} and Chadwick and Sneddon \citep{ChadwickSneddon1958,Chadwick1960}, and it inherits the infinite thermal signal speed. The extensions that remove it are called \emph{generalized} or \emph{non-classical} thermoelasticity \citep{Chandrasekharaiah1986,Chandrasekharaiah1998,HetnarskiIgnaczak1999,HetnarskiIgnaczak2000,IgnaczakOstoja2010}; they now form one of the largest strands of continuum mechanics, with a very large body of theoretical work. Two further developments widened the field. Devices exploit the coupling of deformation and temperature to additional physical fields, above all electric polarization (piezoelectricity and pyroelectricity), but also magnetic induction, diffusing species, fluid-filled pores and microstructural kinematics \citep{Mindlin1974Thermopiezo,Nowacki1978,Eringen1999,Iesan2004}. At small scales, separately, the local stress--strain relation itself fails to describe dispersion, and nonlocal and gradient enrichments are used instead \citep{Eringen1972,Eringen1983,Aifantis1992,Askes2011}.

\subsection{Why another review}
Good reviews of generalized thermoelasticity already exist \citep{Chandrasekharaiah1986,Chandrasekharaiah1998,HetnarskiIgnaczak1999,HetnarskiIgnaczak2000,Shakeriaski2021,Huang2025}. What they do not supply is a treatment for a reader who wants to \emph{use} the theories on wave problems and to compare them reliably. Three observations fixed the organisation of this review.
\begin{enumerate}
\item \textbf{Proliferation without a common language.} The number of named heat-conduction models for thermoelastic waves has grown very large over three decades, and many of them are combinations (two-temperature fractional Lord--Shulman, three-phase-lag with memory, nonlocal Moore--Gibson--Thompson, and so on). Each paper brings its own constants and notation. Section~\ref{sec:framework-link} introduces one heat-flux function $G(s)$ in the Laplace domain that carries every model of the three-lag-plus-memory class (Green--Lindsay and Guyer--Krumhansl do not pass through it, and we treat them separately). A new combination within the class is then a new choice of $G$.
\item \textbf{Unreliable folklore about limits and stability.} Papers repeat several statements without re-deriving them: which parameter limits reduce one model to another, which parameter ranges are stable, how the wave branches should be labelled. We re-derive these statements and check them numerically. One widely repeated claim, that dual-phase-lag thermoelasticity needs $\tau_T\ge\tau_q/2$ to be stable, depends on the order of the Taylor expansion that defines the model (Section~\ref{sec:lags}).
\item \textbf{Figures that cannot be compared.} Published dispersion and response curves use different nondimensionalisations, parameter sets and branch conventions, so that the visual impression of ``the effect of the fractional order'' or ``the effect of nonlocality'' is hard to calibrate. Each figure here states its own scaling (frequency in units of $\omega_*$, or lengths in units of a thermal wavelength), and the material-independent groups that control it are listed in the nomenclature.
\end{enumerate}

\subsection{Contributions and organisation}
No new experiments are reported. The contributions are:
\begin{itemize}
\item a unified heat-flux template, a table of \emph{exact} reductions and singular limits, and a taxonomy that places the main models in relation to one another (Sections~\ref{sec:framework-link}--\ref{sec:memory});
\item a coupled dispersion polynomial valid for the whole template, verified symbolically, with comparative dispersion, damping and speed results (Sections~\ref{sec:classical}--\ref{sec:waves3d});
\item benchmark transient problems, the heat pulse and the thermal shock of a half-space, each with a Laplace-domain closed form and a numerical inversion (Section~\ref{sec:transient});
\item thermoelastic Rayleigh waves, reflection and interface energy partition, each with limit checks (Sections~\ref{sec:surface}--\ref{sec:interfaces});
\item size effects set against lattice dispersion, and the principal multiphysics couplings (Sections~\ref{sec:size}--\ref{sec:multi});
\item a kinetic basis for the Cattaneo and Guyer--Krumhansl laws, with the dispersion of the latter (\cref{sec:kinetic}), and a quantitative test of how well Taylor-truncated lag laws represent the delay law they approximate (\cref{sec:delayfid});
\item flexural thermoelastic damping, including a plate (Lamb-wave) computation validated against the Rayleigh--Lamb equation and the Lifshitz--Roukes solution, resonant modes of a bar, anisotropy in cubic crystals, free-surface reflection and piezoelectric surface waves (\cref{sec:flexted,sec:aniso,sec:reflection,sec:lamb,sec:bar,sec:bg});
\item a discussion of solution methods with finite-difference and inversion benchmarks, experimental evidence and parameter identification, worked examples, application regimes, reporting standards, a comparative synthesis with a pitfalls checklist and a decision guide, and a research agenda (\cref{sec:methods,sec:schemes,sec:applications,sec:scenarios,sec:examples,sec:experiments,sec:synthesis,sec:reporting,sec:open}).
\end{itemize}
Section~\ref{sec:history} gives the historical and literature background, and Section~\ref{sec:foundations} derives the thermomechanical foundations the remaining sections rely on. A reader whose question is simply which model to choose can start at Section~\ref{sec:synthesis}.

\subsection{Scope and conventions}\label{sec:conventions}
We treat \emph{linear} wave propagation in homogeneous solids about a uniform reference state at temperature $T_0$ unless stated otherwise. Nonlinear thermoelasticity, large deformation, radiative transport and chemically reacting media lie outside our scope. The following notation is used throughout.

The Nomenclature table in the front matter collects the principal symbols (\cref{tab:nomenclature}); indices follow the Einstein summation convention.
The strain tensor always carries two indices ($\varepsilon_{ij}$) so that it cannot be confused with the scalar coupling $\varepsilon$. Plane waves are $\exp[\ii(kx-\omega t)]$ with $\Imag k\ge0$ for spatial decay in the propagation direction; a ``branch'' of the dispersion relation is a root $k(\omega)$ followed by continuity in $\omega$.

\subsection{Reproducibility}
Every figure in this review that displays numerical results comes from Python scripts that accompany the manuscript. The identities used repeatedly (the coupled dispersion polynomial, the Debye form of weak-coupling damping, the Green--Naghdi type-II wave speeds, the reductions of Table~\ref{tab:param}) are checked either symbolically or against independent numerical roots in those scripts (Appendix~\ref{app:repro}). Where a quantity depends on a choice of model variant, the form of a memory-dependent derivative for instance, the choice is stated next to the result.

Table~\cref{tab:contributions} collects the contributions a reader can cite this review for. Each row names the place where the contribution is delivered, so that the claim can be checked rather than believed, and the evidence column names only what this paper actually contains.

\begin{table}[H]\centering\small
\caption{What this review contributes, where it is delivered, and what supports it.}
\label{tab:contributions}
\begin{tabularx}{\textwidth}{@{}p{0.5cm}Y p{2.5cm}Y@{}}
\toprule
\textbf{\#} & \textbf{Contribution} & \textbf{Delivered in} & \textbf{Evidence in this paper} \\
\midrule
C1 & One heat-law template and its stepwise reductions, so that LS, GL, GN-II/III, DPL, TPL, MGT, MDD
  and fractional models are compared as members of one family rather than as competing names. &
  \cref{sec:foundations}, \cref{sec:hyperbolic}, \cref{sec:lags}, \cref{sec:memory} &
  \eqref{eq:template} and the reduction table; reductions checked in \path{check_reductions.py} \\
C2 & A positive-definite availability functional with its flux, the dissipation identity it satisfies,
  and the explicit statement that it is not the total-energy flux --- together with the exact Eulerian
  total-energy identity and its first-order reduction. & \cref{sec:foundations}, \cref{sec:interfaces} &
  \eqref{eq:energy}, \eqref{eq:P}, \eqref{eq:exactflux}; net-flux balance at welded interfaces (\cref{sec:interfaces}) \\
C3 & A family-comparative synthesis with a decision guide: what each theory predicts, what it costs, and
  under which conditions two families may legitimately be compared. & \cref{sec:synthesis} &
  comparison tables in \cref{sec:synthesis}; one worked example per family in \cref{sec:examples} \\
C4 & An original five-axis taxonomy that places any candidate model, hybrids included, on one set of
  coordinates, with an explicit dissimilarity that says when two studies are comparable at all. &
  \cref{sec:taxonomy} & axis definitions, the dissimilarity \eqref{eq:taxdissim}, and the assignments in
  \cref{tab:taxonomy-axes} \\
C5 & A per-equation verification register: every registered equation carries a status from a finite
  vocabulary, so that derived, primary-source-checked, unverified and open are distinguishable. &
  \cref{sec:eqregister} & \cref{tab:eqregister}; entries at PRIMARY-SOURCE-CHECKED carry the printed page and equation \\
C6 & Fourteen executable benchmarks with reference solutions and stated tolerances, each a code
  verification rather than a physical validation. & \cref{app:registry}, \cref{app:repro} &
  \cref{app:registry} (14 rows), the scripts of \cref{app:repro}, and the parameter register
  \cref{app:params} \\
C7 & A stated evidence protocol and a set of reporting standards, so that a reader can discount each
  statement by the rule it fails and a future author can make a result comparable. &
  \cref{sec:evidence_protocol}, \cref{sec:reporting} & the six rules, and the reporting checklist \\
\bottomrule
\end{tabularx}
\end{table}

C5 needs one caution, because it is the row a reader is most likely to over-read. A register entry records the status of a check, not the importance or the validity of an equation for a given material. Mechanical correctness does not make a constitutive assumption true.

\subsection{Evidence protocol and how to discount this review}\label{sec:evidence_protocol}

A review is only as useful as the reader's ability to tell what it has checked from what it has merely
repeated. We state the protocol here, so that any statement can be discounted by the rule it fails.

\begin{enumerate}
\item \textbf{Derivation and attribution are separated.} A statement derived here is presented as a
derivation and carries an equation identifier; a statement taken from a source is presented as
attribution even where the source's formulation is reproduced exactly. We never merge the two into one
unmarked sentence. The register in \cref{sec:eqregister} records which is which.
\item \textbf{Every registered equation carries a status} from a fixed vocabulary, and a status records
the check performed, not a rank of importance (\cref{tab:eqregister}).
\item \textbf{A secondary source is never promoted to a primary one.} Where the primary text was not
available the text says so and the register records the weaker status. The absence of the primary text
is not evidence for the claim.
\item \textbf{Bibliographic identifiers are verified before they are printed.} A DOI printed in the
reference list has been checked against a bibliographic record for the same work --- author,
title, journal, volume, pages --- to the level recorded in the reference register (R0--R3), and not
inferred from a citation in another paper. Works that predate DOI assignment say so in the entry instead of carrying a plausible-looking identifier.
\item \textbf{Numbers that can be recomputed are given with what is needed to recompute them.} A benchmark
value is presented with its specification, reference solution and tolerance
(\cref{app:registry}); a value quoted from the literature is presented as quoted, with its conditions,
and not as an independent confirmation.
\item \textbf{Errors are recorded without personal characterisation.} Where a material error is
conclusively established it is recorded plainly; where the evidence establishes a difference but not an
error, it is recorded as a difference. We make no allegation against a publication, and we present no model-specific
discrepancy as a mathematical contradiction without demonstration.
\end{enumerate}

The protocol has limits, and they are part of it. Not every cited work was read page by page; coverage of
the literature is deliberate rather than exhaustive; and the register documents checks performed for this
version, not checks that exist somewhere in principle. For these reasons the register reports coverage
instead of asserting completeness.

\begin{summarybox}[Summary]
This review organises the generalized thermoelasticity literature around one heat-flux template with three lags and one memory term. For linear homogeneous media and rational (truncated-lag) heat laws we derive a dispersion polynomial and check it symbolically against the field equations; exact delays, fractional powers and memory kernels give transcendental or multivalued relations and are treated separately. We add verified benchmarks. Scope: linear waves about a uniform reference state, isotropic and anisotropic solids, plane-wave and guided-wave problems. Contributions: correction to the dual-phase-lag (DPL) stability condition, GN-III front-speed clarification, net-flux-vs-mode-flux check at interfaces, reporting standards. Not covered: nonlinear effects, radiative/reactive heat transfer, full bounded-domain well-posedness literature.
\end{summarybox}
